\section{Singular quadrature, conditioning, and fast application}
\label{sec:em-numerics}

A mesh-free-first integral formulation replaces global volume meshing with
careful treatment of singular kernels, compact local bases, and fast
long-range interactions. These numerical rules are part of the REFERENCE
teacher definition.

\subsection{Interaction classes}

Every pair of source/test supports is classified as:
\begin{enumerate}
\item identical/self;
\item topologically adjacent or overlapping;
\item geometrically near but disjoint;
\item far.
\end{enumerate}
Each class has a separate quadrature policy. One generic Gaussian rule is not
permitted for all four regimes.

\subsection{Self and near-singular terms}

For kernels with $1/R$ or differentiated-$1/R$ singularity, the self block is
split into
\[
  K=K_{\rm sing}+K_{\rm smooth}.
\]
$K_{\rm sing}$ is integrated analytically when possible, or by a
singularity-removing coordinate transformation on the reference
superellipse/surface patch. The smooth remainder is integrated with adaptive
high-order quadrature.

Near interactions use distance-aware refinement or specialized nearly singular
quadrature. A required local estimator is
\[
 \eta_{\rm quad}
 =
 \frac{\norm{K^{(p+1)}-K^{(p)}}}
      {\max(\norm{K^{(p+1)}},\epsilon_{\rm scale})}.
\]
The teacher increases quadrature order or local subdivision until the declared
quadrature budget is met.

\subsection{No fine-minus-coarse self correction}

A conductor self block belongs to the same continuous operator as every other
interaction. It is converged by basis/quadrature enrichment of that block. The
production theory does not define a coarse global self term plus an empirical
fine-minus-coarse correction.

\subsection{Far-field acceleration}

Far interactions are evaluated using one of:
\[
  \text{FMM},\qquad
  \mathcal H/\mathcal H^2\text{-matrices},\qquad
  \text{kernel interpolation / butterfly methods}.
\]
The acceleration tolerance $\epsilon_{\rm far}$ is separate from quadrature and
linear-solve tolerances. The REFERENCE residual uses the same physical operator
with a stricter or independently checked acceleration policy.

\subsection{Block preconditioning}

The preferred preconditioner reflects physics:
\[
 P^{-1}
 \approx
 \operatorname{blockdiag}
 \left(
  A_{\rm cond,self}^{-1},
  A_{\rm surf,self}^{-1},
  A_{\rm contrast,local}^{-1}
 \right)
\]
plus optional coarse/global low-rank coupling. Conductor self inverses may be
cached in intrinsic coordinates and reused under rigid motion in homogeneous
backgrounds.

For current--charge saddle systems, scaling is chosen so current and charge
blocks have comparable nondimensional magnitude. Schur-complement or
constraint preconditioners may be used, but the final residual is always
measured in the original physical block system.

\subsection{Residual norm}

The algebraic residual
\[
 r=b-Ax
\]
is scaled by a block Riesz map $W$:
\[
 \eta_{\rm alg}
 =
 \frac{\sqrt{r\adj W^{-1}r}}
      {\sqrt{b\adj W^{-1}b}}.
\]
$W$ is constructed from physically meaningful current, charge, surface-current,
and contrast-current energy scales. This avoids declaring convergence merely
because one block has numerically small units.

\subsection{Convergence hierarchy}

A REFERENCE electromagnetic sample is accepted only after independently
checking:
\[
 \epsilon_{\rm basis},\quad
 \epsilon_{\rm quad},\quad
 \epsilon_{\rm far},\quad
 \epsilon_{\rm solve}.
\]
Refining one knob cannot compensate for failure of another. In particular, a
small Krylov residual does not certify an under-resolved conductor
cross-section basis.

\subsection{Complexity target}

Let $N$ denote total integral unknowns. Direct dense storage and
$O(N^2)$ application are not production targets. For geometrically regular
scenes the intended REFERENCE operator application is approximately
\[
  O(N\log N)\quad\text{or}\quad O(N)
\]
up to accuracy-dependent constants, with memory of comparable order. The FAST
neural path is independent of $N$.
